% New theory for the paper (20 October 2026): three theorems and one conjecture, written as
% candidates for Section 4 (T1), Section 2/Appendix (T2) and Section 5 (T3).
%
% Status. The statements and proofs below were drafted in this session: for each theorem,
% independent solver runs (pure reasoning, no computation) produced complete proofs, and the
% cleaned proof was then attacked by fresh-context verifiers (general adversarial, one-line-lemma
% extraction, specification check, consecutive pass), 4 agreements to confirm, 2 dissents to
% refute. The vote record is in docs/review/venue_acceptance_plan_2026-10-20.md, Section 9.
% These are drafts for the author to check line by line before anything enters main.tex;
% remarks marked [not voted] were derived by a single solver and not put to the vote.
%
% Notation follows paper/main.tex, Appendix B. References "Lemma B.x", "Theorem 4.x" are to the
% numbering of the build of 20 October 2026 (main at 222d129).
%
% Compile: pdflatex new_theorems_2026-10-20.tex

\documentclass[11pt]{article}
\usepackage[margin=1in]{geometry}
\usepackage{amsmath,amssymb,amsthm}
\usepackage{enumitem}
\usepackage[colorlinks=true,linkcolor=blue]{hyperref}

\newtheorem{theorem}{Theorem}
\newtheorem{lemma}[theorem]{Lemma}
\newtheorem{proposition}[theorem]{Proposition}
\newtheorem{corollary}[theorem]{Corollary}
\newtheorem{conjecture}[theorem]{Conjecture}
\theoremstyle{definition}
\newtheorem{definition}[theorem]{Definition}
\theoremstyle{remark}
\newtheorem{remark}[theorem]{Remark}

\newcommand{\R}{\mathbb R}
\newcommand{\E}{\mathbb E}
\newcommand{\N}{\mathcal N}
\newcommand{\Ppos}{\mathcal P_{>0}}
\newcommand{\nn}{\mathrm{nn}}
\newcommand{\eps}{\varepsilon}
\DeclareMathOperator{\tr}{tr}
\DeclareMathOperator{\Cov}{Cov}
\DeclareMathOperator{\Var}{Var}
\DeclareMathOperator{\aff}{aff}
\DeclareMathOperator{\conv}{conv}
\DeclareMathOperator{\dist}{dist}

\title{New results for ``Atomic Endpoints of Averaged and Guided\\ Empirical Conditional Flows''\\[2pt]
\large Drafts of 20 October 2026, for the author to check}
\date{}

\begin{document}
\maketitle

\paragraph{Summary.}
\begin{enumerate}[label=T\arabic*.,leftmargin=2.5em]
\item \emph{Dominance theorem} (Section~\ref{sec:t1}). A combination $\sum_kc_kv_{p_k}$ of exact prior
  fields with $\sum_kc_k=1$, $K<\infty$, keeps the training atoms for every data set and all positive
  priors whenever its coefficients are convex or one of them dominates,
  $c_j\ge\sum_{k\ne j}|c_k|+\gamma$, uniformly along a schedule of bounded variation. Conversely, every
  constant non-dominant coefficient vector with a negative entry leaves the atoms for some data in
  $\R^2$. It contains Theorem~4.6 and the finite case of Theorem~4.1, explains Theorem~4.7 as a
  non-dominant example, and covers guidance schedules that change sign, which the paper lists as
  open. Theorem~4.1 stays for general mixtures $\Pi$ (Proposition~4.2 needs a continuum of priors,
  where the constants of T1 degenerate).
\item \emph{Schedule invariance} (Section~\ref{sec:t2}). For every Gaussian path
  $X_t=\alpha_tX_1+\sigma_tX_0$ the exact combined flow is, in $y=x/\alpha_t$ and
  $\lambda=\log(\alpha_t/\sigma_t)$, the same equation $dy/d\lambda=\bar m-y$. For paths that start from
  pure noise ($\alpha_0=0$, e.g.\ rectified flow or a variance-preserving path run to $\alpha=0$) the
  endpoint map $x_0\mapsto X_1$ is the same as for the linear path, so all endpoint statements of the
  paper transfer. Samplers that start at a finite noise level (variance-exploding or EDM with finite
  $\sigma_{\max}$, discrete DDPM schedules with $\alpha_T>0$) start the same equation at a finite
  $\lambda_0$; for them only the almost-sure and positive-probability statements transfer.
\item \emph{Off-support sinks survive learned errors} (Section~\ref{sec:t3}). At a non-atom sink of
  Theorem~4.9, a learned field whose scaled error $(1-t)|\Delta|$ stays below a fixed threshold
  still ends within $\limsup(1-t)|P_{L^\perp}\Delta|$ of the non-atom point, with positive
  probability. Along the tie directions the error is suppressed by the factor $1/\beta$.
\item \emph{Conjecture} (Section~\ref{sec:t4}). Stochastic samplers with non-vanishing churn remove
  the off-support endpoints.
\end{enumerate}

\section*{Setting and facts used}

Atoms $a_1,\dots,a_N\in\R^d$ are distinct, $A=\{a_j\}$, $R=\max_j|a_j|$, $\sigma=\min_{i\ne l}|a_i-a_l|$,
$D=\max_{i,l}|a_i-a_l|$. $\Ppos$ is the set of probability vectors on $\{1,\dots,N\}$ with positive
entries. For $t\in(0,1)$: $\eps=1-t$, $\tau=-\log\eps$, $\beta=t/\eps^2=e^{2\tau}-e^\tau$,
$b_j=x\cdot a_j-\frac t2|a_j|^2$, $q_{p,j}=p_je^{\beta b_j}/\sum_lp_le^{\beta b_l}$, $m_p=\sum_jq_{p,j}a_j$,
$v_p=(m_p-x)/\eps$. For $c\in\R^K$ with $\sum_kc_k=1$, $v_c=\sum_kc_kv_{p_k}=(\bar m-x)/\eps$ with
$\bar m=\sum_kc_km_{p_k}$; in the time $\tau$, $\dot x=\bar m-x$. $\nn(z)$ is the set of atoms nearest to
$z$ and $\mathcal C'=\{z:z\in\aff\nn(z)\}$. $\pi_0=\N(0,I_d)$.

We use, from Appendix~B of the paper: Lemma~B.1 (basic facts: $D_xm_p=\beta\Cov_{q_p}(a)$,
$\tr\Cov_q(a)=\sum_jq_j|a_j-y|^2-|m-y|^2$ for every $y$, concentration on $\nn(z)$); the paragraph on
solutions (Carath\'eodory solutions, $C^1$ flow diffeomorphisms, Liouville); Lemma~B.3 (volume
expansion makes non-atom basins null); Lemma~B.4 (signed combinations: for bounded measurable $c$
with $\sum_kc_k=1$, of bounded variation on $[t_0,1)$, every trajectory is bounded, has
$\int_{\tau_0}^\infty|\dot x|^2d\tau<\infty$ and converges to a point of $\mathcal C'$); and
$\operatorname{div}_x(\bar m-x)=\beta\Theta-d$ with $\Theta=\sum_kc_k\tr\Cov_{q_{p_k}}(a)$. For T1(B)
and T3 we use Theorem~4.9(a) and its blow-up coordinates (equation (11), Lemma~B.5).

% ==========================================================================================
\section{T1: the dominance theorem}
\label{sec:t1}

\begin{definition}
Let $c\in\R^K$ with $\sum_kc_k=1$ and $\|c\|_1=\sum_k|c_k|$. Call $c$ \emph{convex} if all
$c_k\ge0$, and \emph{$\gamma$-dominant} ($\gamma>0$) if some index $j$ has
$c_j\ge\sum_{k\ne j}|c_k|+\gamma$; call it \emph{non-dominant} if $\max_kc_k<\|c\|_1/2$. A dominant
coefficient is positive, since $\sum_kc_k=1>0$.
\end{definition}

\begin{theorem}[Dominance theorem]\label{thm:dominance}
\begin{enumerate}[label=(\Alph*),leftmargin=2em]
\item Let $p_1,\dots,p_K\in\Ppos$ and let $c:[0,1)\to\R^K$ be bounded and measurable, with
  $\sum_kc_k(t)=1$, of bounded variation on $[t_0,1)$. Suppose there is $\gamma>0$ such that for
  every $t\in[t_0,1)$, $c(t)$ is convex or $\gamma$-dominant (which one, and which index dominates,
  may change with $t$). Then the initial points whose trajectory of $v_c$ has a non-atom limit
  form a Lebesgue-null set, and under every absolutely continuous source law the endpoint is a
  training atom almost surely.
\item Let $d\ge2$ and let $c\in\R^K$ be constant with $\sum_kc_k=1$, $\min_kc_k<0$ and
  $\max_kc_k<\|c\|_1/2$. Then there are distinct atoms in $\R^d$ and $p_1,\dots,p_K\in\Ppos$ such
  that the endpoint law of $v_c$ from $\pi_0$ puts positive mass on a non-atom point.
\end{enumerate}
\end{theorem}

\paragraph{Notation for the proof.} $q_k=q_{p_k}(x,t)$, $m_k=m_{p_k}$, $T_k=\tr\Cov_{q_k}(a)\in[0,R^2]$,
$C=\sup_t\|c(t)\|_1$, $\rho=\max_{k,l,i}p_{k,i}/p_{l,i}\ge1$.

\begin{lemma}[Three elementary facts]\label{lem:t1facts}
\begin{enumerate}[label=(\roman*),leftmargin=2em]
\item $\tr\Cov_q(a)\ge\frac{\sigma^2}2(1-\max_iq_i)$ for every probability vector $q$.
\item $q_{l,i}\le\rho^2q_{k,i}$ for all $k,l,i$ and all $(x,t)$.
\item For $z\in\R^d$ with $r=\dist(z,A)$, $S=\nn(z)$, $u_k=m_k-z$ and
  $A_k=\sum_iq_{k,i}|a_i-z|^2$: $T_k=A_k-|u_k|^2$, $|u_k|^2\le A_k$, and
  $A_k=r^2+\sum_{a_i\notin S}q_{k,i}(|a_i-z|^2-r^2)\ge r^2$.
\end{enumerate}
\end{lemma}
\begin{proof}
(i) Expanding around the mean, $\sum_{i,l}q_iq_l|a_i-a_l|^2=2\tr\Cov_q(a)$, so
$\tr\Cov_q(a)\ge\frac{\sigma^2}2\sum_{i\ne l}q_iq_l=\frac{\sigma^2}2(1-\sum_iq_i^2)\ge\frac{\sigma^2}2(1-\max_iq_i)$.
(ii) With $Z_k=\sum_ip_{k,i}e^{\beta b_i}$, $q_{l,i}/q_{k,i}=(p_{l,i}/p_{k,i})(Z_k/Z_l)$ and
$Z_k/Z_l\le\max_ip_{k,i}/p_{l,i}\le\rho$. (iii) Lemma~B.1 with $y=z$, and Jensen.
\end{proof}

\begin{lemma}[Convex times]\label{lem:t1convex}
Let $z\notin A$, $r=\dist(z,A)$. If $c$ is convex and $|\bar m-z|\le r/2$, then
$\Theta\ge\theta_c:=\sigma^2r/(4KD\rho^2)$.
\end{lemma}
\begin{proof}
Pick $k_0$ with $c_{k_0}\ge1/K$, let $i^*$ maximise $q_{k_0,i}$ and $\delta=1-q_{k_0,i^*}$. As all
$T_k\ge0$, Lemma~\ref{lem:t1facts}(i) gives $\Theta\ge T_{k_0}/K\ge\sigma^2\delta/(2K)$. By
Lemma~\ref{lem:t1facts}(ii), $\sum_{i\ne i^*}q_{l,i}\le\rho^2\delta$ for every $l$, so
$|m_l-a_{i^*}|\le D\rho^2\delta$ and, by convexity, $|\bar m-a_{i^*}|\le D\rho^2\delta$. Then
$r\le|z-a_{i^*}|\le r/2+D\rho^2\delta$, so $\delta\ge r/(2D\rho^2)$.
\end{proof}

\begin{lemma}[Dominant times]\label{lem:t1dominant}
Let $z\notin A$, $r=\dist(z,A)$, $\Delta=\bar m-z$, and
$e=\max_k\sum_{a_i\notin\nn(z)}q_{k,i}(|a_i-z|^2-r^2)\ge0$. If $c$ is $\gamma$-dominant, then
\[
\Theta\ \ge\ \gamma r^2-(C+1)e-2\big(2C|\Delta|\sqrt{r^2+e}+|\Delta|^2\big).
\]
\end{lemma}
\begin{proof}
Let $j$ be the dominant index, $P=\sum_{k\ne j,c_k>0}c_k$, $Q=\sum_{k\ne j,c_k<0}|c_k|$, $W=P+Q\le C$.
From $c_j+P-Q=1$ and $c_j\ge W+\gamma$: $2P\le1-\gamma$ (so $\gamma\le1$), $c_j\ge1-P\ge\frac12$,
$W<c_j$. By
Lemma~\ref{lem:t1facts}(iii), $\Theta=\sum_kc_kA_k-F$ with $F=\sum_kc_k|u_k|^2$, and
$\sum_kc_kA_k\ge r^2-Qe$ because only indices $k\ne j$ can carry negative coefficients and
$A_k-r^2\in[0,e]$. Put $V=\sum_{k\ne j}c_ku_k$. Since $\sum_kc_ku_k=\Delta$, $c_ju_j=\Delta-V$; also
$|V|\le W\sqrt{r^2+e}$ and, by Cauchy--Schwarz, $|V|^2\le W\sum_{k\ne j}|c_k||u_k|^2$. Hence
\[
c_j|u_j|^2\le\frac{(|V|+|\Delta|)^2}{c_j}\le\frac{|V|^2}{c_j}+E,\qquad
E=2\big(2|\Delta|W\sqrt{r^2+e}+|\Delta|^2\big),
\]
using $c_j\ge\frac12$, so $F\le\sum_{k\ne j}\big(W|c_k|/c_j+c_k\big)|u_k|^2+E$. For $c_k<0$ the
coefficient is $|c_k|(W/c_j-1)\le0$; for $c_k>0$ it is at most $2c_k$, and $|u_k|^2\le r^2+e$. So
$F\le2P(r^2+e)+E\le(1-\gamma)(r^2+e)+E$ and $\Theta\ge\gamma r^2-(Q+1-\gamma)e-E$.
\end{proof}

\begin{proof}[Proof of Theorem~\ref{thm:dominance}(A)]
By Lemma~B.4 every trajectory is bounded, has $\int_{\tau_0}^\infty|\dot x|^2<\infty$ and converges to a
point of $\mathcal C'$; $|\bar m|\le CR$. By Lemma~B.3 it suffices to show
$\int_{\tau_0}^n(\beta\Theta-d)\,d\tau\to+\infty$ along every trajectory whose limit $z$ is not an atom.
Fix one, put $r=\dist(z,A)>0$, $\Delta=\bar m-z$, $g=|\bar m-x|=|\dot x|$ a.e., and $e(\tau)$ as in
Lemma~\ref{lem:t1dominant}; $e(\tau)\to0$ by Lemma~B.1(c), and $|\Delta|\le g+|x-z|$.

\emph{Constants.} $\theta_0=\min(\gamma r^2/2,\theta_c)$. Choose $e_0\in(0,r^2]$ with
$(C+1)e_0\le\gamma r^2/4$, then $\eta\in(0,r/4]$ with $2(4\sqrt2\,Cr\eta+4\eta^2)\le\gamma r^2/4$, then
$\tau_1\ge\tau_0$ with $|x-z|\le\eta$ and $e\le e_0$ for $\tau\ge\tau_1$. At every $\tau\ge\tau_1$ with
$g(\tau)\le\eta$ we have $|\Delta|\le2\eta\le r/2$, and Lemma~\ref{lem:t1convex} or
Lemma~\ref{lem:t1dominant} gives $\Theta\ge\theta_0$. Always $\Theta\ge-M$ with $M=CR^2+1$.

\emph{Divergence.} Raising $t_0$ if necessary, $\tau_0>0$. $\beta$ is positive and increasing for
$\tau>0$, and
$\beta(\tau+1)/\beta(\tau)\le\Lambda:=e^2+e(e-1)/(e^{\tau_0}-1)$ for $\tau\ge\tau_0>0$. For integers
$n\ge\lceil\tau_1\rceil$ let $I_n=[n,n+1]$ and $B_n=\{\tau\in I_n:g>\eta\}$; by Chebyshev,
$|B_n|\le\eta^{-2}\int_{I_n}g^2\to0$. Take $n_2\ge\max(\lceil\tau_1\rceil,\tau_0+1)$ with
$|B_n|\le\min\{\frac14,\theta_0/(4M\Lambda^2)\}$ for $n\ge n_2$. Then
\[
\int_{I_n}\beta\Theta\ \ge\ \beta(n)\theta_0(1-|B_n|)-M\beta(n+1)|B_n|
\ \ge\ \beta(n)\Big(\tfrac34\theta_0-\tfrac{\theta_0}{4\Lambda}\Big)\ \ge\ \tfrac12\theta_0\beta(n),
\]
so $J_n=\int_{I_n}(\beta\Theta-d)\ge\frac12\theta_0\beta(n)-d\to\infty$, and
$\int_{\tau_0}^n\operatorname{div}f=\int_{\tau_0}^{n_2}\operatorname{div}f+\sum_{m=n_2}^{n-1}J_m\to+\infty$.
Lemma~B.3 concludes.
\end{proof}

\begin{lemma}[Short resultants in the plane]\label{lem:t1polygon}
Let $\ell_1,\dots,\ell_n>0$ with sum $P$ and maximum $\ell_{\max}$. There are unit vectors
$s_i\in\R^2$ with $|\sum_i\ell_is_i|=\max(0,2\ell_{\max}-P)$.
\end{lemma}
\begin{proof}
If $2\ell_{\max}\ge P$, point the largest vector opposite all the others. Otherwise sort decreasingly
and let $m$ be minimal with $\sum_{i\le m}\ell_i\ge P/2$; the three group lengths
$x=\sum_{i<m}\ell_i$, $\ell_m$ and $y=\sum_{i>m}\ell_i$ are each at most $P/2$, hence satisfy the
triangle inequalities, and three such lengths close to a triangle. Point every vector of a group
along its side of the triangle.
\end{proof}

\begin{lemma}[Prescribed means]\label{lem:t1means}
If $A\subset\R^2$ is finite and $y\in\operatorname{int}\conv A$, some $p\in\Ppos$ (on $A$) has
$\E_p[a]=y$.
\end{lemma}
\begin{proof}
Let $U$ be uniform on $A$ with mean $\bar u$. For small $\alpha>0$,
$y_\alpha=(y-\alpha\bar u)/(1-\alpha)\in\conv A$; write $y_\alpha=\sum_ir_ia_i$ and take
$p=\alpha U+(1-\alpha)r$, which has $p_i\ge\alpha/|A|$ and mean $y$.
\end{proof}

\begin{proof}[Proof of Theorem~\ref{thm:dominance}(B)]
Let $P_\pm=\{k:\pm c_k>0\}$, $c_+=\sum_{P_+}c_k$, $c_-=\sum_{P_-}|c_k|>0$, so $c_+=1+c_-$, and let
$c_j=\max_kc_k$. Non-dominance means $2c_j-c_+<c_-$.

\emph{Directions.} By Lemma~\ref{lem:t1polygon} applied to $(c_k)_{k\in P_+}$ in a $2$-plane
$E\subset\R^d$, there are unit vectors $v_k\in E$ with $w=\sum_{P_+}c_kv_k$ and
$|w|=\max(0,2c_j-c_+)<c_-$. Put $y=w/c_-$, so $|y|<1$.

\emph{Atoms.} Pick $n\ge3$ with $\cos(\pi/n)>|y|$ and let $A$ consist of the distinct points among
the $v_k$ together with the vertices of a regular $n$-gon on the unit circle of $E$. All atoms have
norm $1$, $A$ contains three non-collinear points, and $\conv A$ contains the disc of radius
$\cos(\pi/n)$, so $0$ and $y$ are interior points of $\conv A$ in $E$. Lemma~\ref{lem:t1means} gives
$p_0,p_y\in\Ppos$ with means $0$ and $y$. Since $A$ is not contained in a line, $\Cov_p(a)\succ0$ on
$E$ for every $p\in\Ppos$; let $\lambda_0>0$ be the smaller of the least eigenvalues on $E$ of
$\Cov_{p_0}$ and $\Cov_{p_y}$.

\emph{Priors.} Let $\delta=\min\{\frac12,c_-\lambda_0/(4c_+)\}$ and set
$p_k=(1-\delta)\mathbf 1_{v_k}+\delta p_0$ for $k\in P_+$ (mean $(1-\delta)v_k$),
$p_k=(1-\delta)p_y+\delta p_0$ for $k\in P_-$ (mean $(1-\delta)y$), and $p_k=p_0$ if $c_k=0$. All lie
in $\Ppos$.

\emph{Verification.} Take $z=0$. Every atom is at distance $1$, so $\nn(0)=A$, $\aff A=E\ni0$, and
$0$ is a non-atom point of $\mathcal C'$ with $S=A$ and $L=E$. At the profile $\xi^\star=0$,
$\pi_k(0)=p_k$ and $M(0)=(1-\delta)(w-c_-y)=0=z$. For $k\in P_+$,
$\tr\Cov_{p_k}=1-(1-\delta)^2\le2\delta$ since all atoms are unit vectors, so
$\Cov_{p_k}\preceq2\delta I$; for $k\in P_-$, the covariance of a mixture dominates the mixture of
covariances, so $\Cov_{p_k}\succeq\lambda_0I$. Hence
$\Gamma(0)\preceq(2\delta c_+-c_-\lambda_0)I\preceq-\frac12c_-\lambda_0I\prec0$ on $L$, and
Theorem~4.9(a) gives positive probability of the endpoint $0$ under $\pi_0$.
\end{proof}

\begin{corollary}[Consequences]\label{cor:dominance}
Let all priors lie in $\Ppos$ (for label-conditioned fields: $h>0$).
\begin{enumerate}[label=(\roman*),leftmargin=2em]
\item \emph{Two fields, any bounded schedule.} For $c(t)=(1+w(t),-w(t))$ with $w$ bounded,
  measurable and of bounded variation on $[t_0,1)$, of any sign, every $c(t)$ is convex
  ($w\in[-1,0]$) or $1$-dominant ($c_1-|c_2|=1$ for $w>0$, $c_2-|c_1|=1$ for $w<-1$). So guidance
  with any bounded schedule of bounded variation, including schedules that change sign
  \cite{jain2025attraction}, keeps the atoms; and $p_u$ need not be uniform.
\item \emph{One positive coefficient.} If exactly one $c_j>0$, then
  $c_j-\sum_{k\ne j}|c_k|=\sum_kc_k=1$. This covers guidance with several negative prompts,
  $(1+\sum_iw_i)v_c-\sum_iw_iv_{n_i}$, and guidance against a label-averaged field,
  $(1+w)v_p-w\sum_k\lambda_kv_{p_k}$.
\item \emph{Guided label averaging.} For $c=((1+w)\lambda_1,\dots,(1+w)\lambda_K,-w)$, $w>0$, $c$ is
  dominant iff $\lambda_{\max}>(1+2w)/(2(1+w))$, i.e.\ the averaging keeps more than
  $1-1/(2(1+w))$ of its weight on one label; then $\gamma=(1+w)(2\lambda_{\max}-1)-w$. For
  $\lambda=(\frac12,\frac12)$ it is non-dominant for every $w>0$, consistent with Theorem~4.7.
\item \emph{Composable conjunctions.} For $c=(1-\sum_iw_i,w_1,\dots,w_n)$ with $w_i>0$ and
  $W=\sum_iw_i>1$ (the AND operator of composable diffusion), $c$ is dominant iff
  $W-w_{\max}<\frac12$. Two concepts of equal weight $w>\frac12$ are non-dominant.
\end{enumerate}
\end{corollary}

\begin{remark}[The constant, and dimension one; not voted]\label{rem:kappa}
With $\kappa(c)=\sup\{\sum_kc_k|u_k|^2:u_k\in\R^m,\ |u_k|\le1,\ \sum_kc_ku_k=0\}$, $m\ge2$, $A$ the
largest coefficient, $B$ the other positive mass and $C$ the negative mass, one solver obtained
$\kappa=B/(1-B)$ if $B\le\frac12$ and $\kappa=1+C-(A-B)_+^2/C$ if $B\ge\frac12$; so $\kappa<1$ exactly
under dominance and $\kappa>1$ for non-dominant $c$ with a negative entry. In dimension one the
criterion changes: the solver shows that $c=(\frac7{20},\frac7{20},\frac7{20},-\frac1{20})$ is
non-dominant but keeps the atoms for all data on a line (the scalar version of $\kappa$ is a
signed-sum problem), while every non-dominant $c$ with at most two positive entries leaves already
on two atoms. Part (B) therefore needs tie dimension at least two in general. These two claims
should be checked before they are used.
\end{remark}

\begin{remark}[Open boundary]
The boundary $\max_kc_k=\|c\|_1/2$ with a negative entry, and schedules whose dominance margin tends
to zero, are not covered.
\end{remark}

% ==========================================================================================
\section{T2: schedule invariance}
\label{sec:t2}

A \emph{Gaussian path} is $(\alpha,\sigma)$ on $[0,1]$, continuous on $[0,1]$ and $C^1$ on $[0,1)$, with
$\alpha_0=0$, $\sigma_0=1$, $\alpha_1=1$, $\sigma_1=0$, $\alpha_t>0$ on $(0,1]$, $\sigma_t>0$ on $[0,1)$, and
$\lambda_t=\log(\alpha_t/\sigma_t)$ of class $C^1$ on $(0,1)$ with $\lambda'>0$. For $p\in\Ppos$ let
$I\sim p$, $X_1=a_I$, $X_0\sim\N(0,I_d)$ independent, $X_t=\alpha_tX_1+\sigma_tX_0$, and
$v^{\alpha,\sigma}_p(x,t)=\E[\alpha_t'X_1+\sigma_t'X_0\mid X_t=x]$ (continuous Bayes version). For
$c$ bounded measurable with $\sum_kc_k=1$, $V_c=\sum_kc_k(t)v^{\alpha,\sigma}_{p_k}$. Let
$\tilde m_p(y,\lambda)=\sum_j\pi_{p,j}a_j$ with $\pi_{p,j}\propto p_j\exp(e^{2\lambda}(y\cdot a_j-|a_j|^2/2))$.

\begin{theorem}[Schedule invariance]\label{thm:schedule-invariance}
\begin{enumerate}[label=(\alph*),leftmargin=2em]
\item $v^{\alpha,\sigma}_p=\frac{\sigma'}\sigma x+\big(\alpha'-\frac{\sigma'\alpha}\sigma\big)m_p$, where
  $m_p(x,t)=\tilde m_p(x/\alpha_t,\lambda_t)$, and $\alpha'-\sigma'\alpha/\sigma=\alpha\lambda'$.
\item $x(t)$ is a Carath\'eodory solution of $\dot x=V_c$ on an interval $J\subset(0,1)$ iff
  $Y(\lambda)=x(t)/\alpha_t$, $\lambda=\lambda_t$, solves
  \[
  \frac{dY}{d\lambda}=\sum_k\tilde c_k(\lambda)\,\tilde m_{p_k}(Y,\lambda)-Y,\qquad\tilde c(\lambda_t)=c(t),
  \]
  on $\lambda(J)$. For the linear path this is the paper's equation $\dot x=\bar m-x$ in the time
  $\tau$, under $Y=x/t$ and $\lambda=\log(t/(1-t))=\log(e^\tau-1)$.
\item For two Gaussian paths and schedules with the same $\tilde c$, the solutions from the same
  $x_0$ at $t=0$ exist and are unique on $[0,1)$, and $\lim_{t\uparrow1}x(t)$ exists for one iff
  for the other, with the same value. Explicitly,
  $x(t)=(\alpha_t+\sigma_t)\,x_{\rm lin}\big(\alpha_t/(\alpha_t+\sigma_t)\big)$.
\item Consequently every statement about the endpoint law under $\N(0,I_d)$ at $t=0$ that holds for
  the linear path holds verbatim for every Gaussian path, with schedules read as functions of
  $\lambda$. Combinations with $\sum_kc_k=1$ of data predictions, $\eps$-predictions, scores or
  velocities at the same $(x,t)$ define the same combined field.
\end{enumerate}
\end{theorem}

\begin{proof}
\emph{(a)} Given $I=j$ and $X_t=x$, $X_0=(x-\alpha_ta_j)/\sigma_t$, so
$\alpha'X_1+\sigma'X_0=\frac{\sigma'}\sigma X_t+\big(\alpha'-\frac{\sigma'\alpha}\sigma\big)X_1$ pointwise, and
Bayes' rule gives the formula with $m_p=\sum_jq_ja_j$, $q_j\propto p_j\exp(-|x-\alpha a_j|^2/2\sigma^2)$;
the right side is continuous, hence the Bayes version. Expanding the square, the $j$-independent term
cancels and $(\alpha/\sigma^2)x\cdot a_j-(\alpha^2/2\sigma^2)|a_j|^2=e^{2\lambda}(y\cdot a_j-|a_j|^2/2)$ with
$y=x/\alpha$. Finally $\lambda'=\alpha'/\alpha-\sigma'/\sigma$ gives $\alpha\lambda'=\alpha'-\sigma'\alpha/\sigma$ and
$\sigma'/\sigma-\alpha'/\alpha=-\lambda'$.

\emph{(b)} Using $\sum_kc_k=1$, $V_c=\frac{\sigma'}\sigma x+\alpha\lambda'\,M$ with
$M=\sum_k\tilde c_k(\lambda_t)\tilde m_{p_k}(x/\alpha_t,\lambda_t)$. For $y=x/\alpha$, which is absolutely
continuous on compacts, a.e.\ $y'=(\frac{\sigma'}\sigma-\frac{\alpha'}\alpha)y+\lambda'M=\lambda'(M-y)$. The
substitution rule for the $C^1$ increasing bijection $\lambda$ (valid for integrands bounded on compacts;
it also makes $\tilde c=c\circ\lambda^{-1}$ measurable, since $\lambda^{-1}$ is locally Lipschitz) turns
this into the integral form of $dY/d\lambda=\tilde M-Y$, and conversely; from $y'=\lambda'(M-y)$,
$x'=\alpha'y+\alpha\lambda'(M-y)=\frac{\sigma'}\sigma x+\alpha\lambda'M$. For the linear path,
$\sigma'/\sigma=-1/(1-t)$, $\alpha\lambda'=1/(1-t)$, $e^{2\lambda}=t^2/(1-t)^2$, and the weights become
$p_j\exp(\beta(x\cdot a_j-\frac t2|a_j|^2))$.

\emph{(c)} \emph{Existence and uniqueness from $t=0$.} In $(\lambda,y)$ the start is singular
($\lambda\to-\infty$, $\lambda'$ may blow up), so we work in $x$. On $[0,T]$, $T<1$, the functions
$\sigma'/\sigma$, $\alpha\lambda'=\alpha'-\sigma'\alpha/\sigma$ and $\theta=\alpha/\sigma^2$ are continuous (with
$\theta(0)=0$), and $D_xm_p=\theta\Cov_q$ with $\|\Cov_q\|\le R^2$. So $V_c$ has linear growth and is
Lipschitz in $x$ with an integrable constant, and a Bielecki-norm contraction for
$x\mapsto x_0+\int_0^tV_c(x(s),s)\,ds$ gives a unique solution on $[0,T]$, hence on $[0,1)$, depending
continuously on $x_0$.

\emph{Same $\lambda$-curve.} Let $Y(\mu)=x(\lambda^{-1}(\mu))/\alpha$ and $\hat Y$ the analogue for the
other path. Both solve the equation of (b) on $\R$ with the same $\tilde c$, and
$e^\mu Y(\mu)=x(t)/\sigma_t\to x_0$ as $\mu\to-\infty$, likewise for $\hat Y$. Let
$W=e^\mu(Y-\hat Y)$. The $-Y$ term cancels the derivative of $e^\mu$, so a.e.
$W'=e^\mu\sum_k\tilde c_k(\tilde m_{p_k}(Y)-\tilde m_{p_k}(\hat Y))$, and since
$D_y\tilde m_p=e^{2\mu}\Cov_\pi$, $|W'|\le k(\mu)|W|$ with $k=CR^2e^{2\mu}$, integrable on
$(-\infty,\mu_1]$. As $W\to0$ at $-\infty$, $|W(\mu)|\le\Phi(\mu):=\int_{-\infty}^\mu k|W|$,
$\Phi'\le k\Phi$, and with $K(\mu)=\int_{-\infty}^\mu k$, $\Phi e^{-K}$ is nonincreasing with limit $0$ at
$-\infty$; so $\Phi\equiv0$ and $Y=\hat Y$. Since $\alpha_t\to1$ and $\lambda_t\to\infty$ as $t\uparrow1$,
$\lim x$ exists iff $\lim_{\mu\to\infty}Y$ exists, with the same value, and the same holds for $\hat x$. The
explicit relation follows from $\lambda_{\rm lin}(s)=\lambda_t$ at $s=\alpha_t/(\alpha_t+\sigma_t)$.

\emph{(d)} The set $D$ of $x_0$ with a limit is Borel (Cauchy criterion over rational times, joint
continuity in $(t,x_0)$) and the endpoint map is Borel on $D$; both are the same for every path, so
$\Pr(X_0\in D,\ X_1\in S)$ is the same for every Borel $S$. Each of data prediction, $\eps$-prediction,
score and velocity is an affine function $a+b\,m_p$ of $m_p$ with $p$-independent coefficients and
$b\ne0$, so sums with $\sum_kc_k=1$ commute with the conversions.
\end{proof}

\begin{remark}
\emph{Why $\sum_kc_k=1$.} With $s=\sum_kc_k$ the $\lambda$-equation gains $(s-1)\frac{d\log\sigma}{d\lambda}y$,
which depends on the path ($-e^\lambda/(1+e^\lambda)$ for the linear path, $-e^{2\lambda}/(1+e^{2\lambda})$ for
$\alpha=\sin\frac{\pi t}2$, $\sigma=\cos\frac{\pi t}2$). \emph{Normalisations.} With $\sigma_0\ne1$ the endpoint
map is $x_0\mapsto\Phi_{\rm lin}(x_0/\sigma_0)$; schedules must be matched in $\lambda$, not in $t$.
\emph{Finite start.} If $\alpha_0>0$ (variance-exploding or EDM with finite $\sigma_{\max}$, discrete
DDPM with $\alpha_T>0$), the sampler starts the $\lambda$-equation at a finite $\lambda_0$ from a Gaussian
that is not the path marginal there, so the endpoint map and the endpoint weights (e.g.\ $p_j$ in
Theorem~3.4) differ from the linear path's. Statements of the form ``the endpoint is an atom almost
surely'' and ``the endpoint equals a given non-atom point with positive probability'' still hold:
the flow maps are diffeomorphisms, Lemma~B.3 applies from any starting time, and the open basins of
Theorem~4.9(a) pull back to nonempty open sets, which a Gaussian charges.
\emph{Novelty.} The reduction to $(x/\alpha,\log\text{SNR})$ is the empirical-target case of the schedule
invariance of probability-flow ODEs used by Karras et al.\ (2022); the contribution here is only
that all endpoint statements of the paper transfer, including guided and averaged fields. The result
concerns deterministic samplers.
\end{remark}

% ==========================================================================================
\section{T3: off-support sinks survive learned errors}
\label{sec:t3}

Fix $z$, $S=\nn(z)$ with $|S|\ge2$ and $z\in\aff S$, $L=\operatorname{span}\{a-z:a\in S\}$, constant $c$ with
$\sum_kc_k=1$ and $p_k\in\Ppos$, and $M$, $G=M-z$, $\Gamma=DG$ as in Theorem~4.9. Blow-up coordinates:
$\xi=\beta P_L(x-tz)$, $\eta=P_{L^\perp}(x-z)$, so $x-z=\xi/\beta-\eps P_Lz+\eta$.

\begin{theorem}[Robust off-support sinks]\label{thm:robust-sink}
Let $\xi^\star$ be a profile with $\Gamma(\xi^\star)\prec0$ on $L$. There are $\rho>0$, $\eta_0>0$ and
$t_1\in(0,1)$ such that the following holds for every $\Delta:\R^d\times[0,1)\to\R^d$ that is measurable
in $t$, $C^1$ in $x$ with locally bounded derivative and of linear growth on each $[0,T]$, $T<1$, and
satisfies
\[
(1-t)\,|\Delta(x,t)|\le\eta_0\quad\text{for }t\in[t_1,1),\ x\in W_t:=\{|\beta P_L(x-tz)-\xi^\star|\le\rho,\
|P_{L^\perp}(x-z)|\le\rho\}.
\]
Let $O$ be the set of $x_0$ whose trajectory of $v_c+\Delta$ lies in $\operatorname{int}W_{t_1}$ at
time $t_1$. Then $O$ is open and nonempty, and for $x_0\in O$ the trajectory stays in $W_t$ for all
$t\in[t_1,1)$ and
\begin{enumerate}[label=(\roman*),leftmargin=2em]
\item $P_L(x_t-z)\to0$;
\item $\limsup_{t\uparrow1}|x_t-z|\le\limsup_{t\uparrow1}(1-t)|P_{L^\perp}\Delta(x_t,t)|\le\eta_0$;
\item if $(1-t)\sup_{W_t}|\Delta|\to0$, then $x_t\to z$.
\end{enumerate}
In particular the perturbed flow stays asymptotically within $\eta_0$ of the non-atom $z$ with
probability $\pi_0(O)>0$.
\end{theorem}

\begin{proof}
\emph{Perturbed blow-up equations.} In the time $\tau$, $\dot x=\bar m-x+\eps\Delta$ a.e. With
$\dot\beta/\beta-1=1/t$ and $1/(t\beta)=\eps^2/t^2$,
\[
\frac{d\xi}{d\tau}=\beta\big[G(\xi)+\mathcal R+\eps P_L\Delta\big],\qquad
\frac{d\eta}{d\tau}=-\eta+\mathcal E+\eps P_{L^\perp}\Delta,
\]
with $\mathcal R=(\eps^2/t^2)\xi+P_L(\bar m-M(\xi))$ and $\mathcal E=P_{L^\perp}(\bar m-z)$ as in equation~(11):
$P_L(\bar m-x-\eps z)=P_L(\bar m-z)-\xi/\beta$, and $P_L(\bar m-z)=G(\xi)+P_L(\bar m-M(\xi))$ because
$M(\xi)-z\in L$ (here $\sum_kc_k=1$ is used).

\emph{Monotonicity.} Let $\lambda>0$ be the least eigenvalue of $-\Gamma(\xi^\star)$ and choose $\rho_0>0$
with $\|\Gamma(\xi)-\Gamma(\xi^\star)\|\le\lambda/2$ for $|\xi-\xi^\star|\le\rho_0$. Integrating $\Gamma$ along
the segment, $\langle h,G(\xi^\star+h)\rangle\le-\frac\lambda2|h|^2$ for $h\in L$, $|h|\le\rho_0$.

\emph{Constants, fixed before $\Delta$.} Apply Lemma~B.5(i) to the compact
$K=\overline B_L(\xi^\star,\rho_0)$ to get $\bar\rho$ and $\bar\tau$. Put $\rho=\min(\rho_0,\bar\rho)$ and
$\eta_0=\rho\min(\lambda/4,1/2)$, and choose $\tau_1\ge\max(\bar\tau,1)$ such that
$|\mathcal R|\le\lambda\rho/8$ and $|\mathcal E|\le\rho/4$ on $\{\xi\in K,|\eta|\le\bar\rho\}$ for
$\tau\ge\tau_1$; $t_1=1-e^{-\tau_1}$. Then $W_t$ lies in the region of Lemma~B.5(i).

\emph{Inequalities.} At a.e.\ $\tau\ge\tau_1$ with $x\in W_\tau$, for $h=\xi-\xi^\star$, $V=|h|^2$, $U=|\eta|^2$:
$V'\le\lambda\beta|h|(\frac34\rho-|h|)$ and $U'\le2|\eta|(\frac34\rho-|\eta|)$.

\emph{Invariance.} If $u$ is absolutely continuous on $[a,b]$ and $u'\le0$ a.e.\ on $\{u>c'\}$, then
$u(b)\le\max(c',u(a))$. Let $\mathcal X=\{\tau\ge\tau_1:x(\tau)\notin W_\tau\}$ and suppose it is nonempty,
with $T=\inf\mathcal X$. Its complement is closed, so $x\in W$ on $[\tau_1,T]$, and the barrier statement with
$c'=9\rho^2/16$ gives $V(T),U(T)<\rho^2$, which persists on $[T,T+\delta)$, contradicting $T=\inf\mathcal X$.
So $x_t\in W_t$ for $t\ge t_1$ and $|\xi|\le|\xi^\star|+\rho$.

\emph{Conclusions.} (i) $|P_L(x_t-z)|\le(|\xi^\star|+\rho)/\beta+\eps|z|\to0$. (ii) By variation of
constants for $\eta$, $\limsup|\eta|\le\limsup(|\mathcal E|+\eps|P_{L^\perp}\Delta|)$, and $|\mathcal E|\to0$
on $W$; combine with (i). (iii) follows from (ii). $O$ is the preimage of the nonempty open set
$\operatorname{int}W_{t_1}$ under the flow homeomorphism from $0$ to $t_1$, so $\pi_0(O)>0$.
\end{proof}

\begin{remark}[Sharpness]
(a) The bound (ii) is attained: for a unit $e\in L^\perp$ and $\Delta=\eta_0e/(1-t)$, $x_t\to z+\eta_0e$.
(b) A tangential error $\Delta=\eta_0e/(1-t)$, $e\in L$, still gives $x_t\to z$: in the tie directions the
restoring stiffness is $\lambda\beta\to\infty$, so the displacement $\xi/\beta$ is $O(\eta_0/(\lambda\beta))$.
If $\aff S=\R^d$, the endpoint is exactly $z$ for every admissible $\Delta$. (c) $\eta_0$ must be of
order $\rho$; and $\pi_0(O)$ has no lower bound uniform in $\Delta$, because $\Delta$ is unrestricted
before $t_1$.
\end{remark}

\begin{remark}[Contrast with an atom]
Near an atom ($L=\{0\}$) the restoring rate is $1$ in $\tau$ in every direction, the same scale as the
forcing $\eps\Delta$, so an endpoint error of order $\limsup(1-t)|\Delta|$ appears in every direction
(Theorem~5.1). At an off-support sink only the component normal to the tie survives, so the
endpoint error is smaller. The admissible error is not larger, however: $\eta_0=\rho\min(\lambda/4,1/2)$
shrinks with the spectral gap $\lambda$ of $-\Gamma(\xi^\star)$, which tends to zero as the guidance
strength approaches the threshold, while an atom tolerates errors of order $\rho/4$
(Proposition~5.2).
\end{remark}

% ==========================================================================================
\section{T4: stochastic samplers (conjecture)}
\label{sec:t4}

\begin{conjecture}\label{conj:sde}
For a stochastic sampler whose noise has constant churn in log-SNR time (for example the reverse
SDE, or EDM's stochastic sampler with $S_{\rm churn}>0$ held fixed), applied to any bounded
combination $\sum_kc_kv_{p_k}$ with $p_k\in\Ppos$, the endpoint is a training atom almost surely.
In particular guided label averaging produces no off-support endpoints under such samplers.
\end{conjecture}

\paragraph{Heuristic.} In the profile coordinate $\xi=\beta P_L(x-tz)$ the noise of size
$\sigma_t\sqrt{d\lambda}$ in $y$ becomes $O(1)$ per unit of the fast time $\sigma$, because the profile
scale is $1/\beta=e^{-2\lambda}$ and $d\sigma=\beta\,d\tau$. The profile then follows a diffusion with drift
$\nabla\Psi$ (Theorem~4.9) and non-vanishing noise over an infinite fast-time horizon, which leaves every
bounded neighbourhood of a sink, while far from the profiles the drift has a definite size.
\emph{What a proof needs:} a transience statement for this nonautonomous diffusion, and convergence
of $x$ to an atom once $\xi$ escapes. \emph{Test:} Euler--Maruyama on the exact field of Theorem~4.7
($B=2$, $w=5w_c$) with churn $\zeta\in\{0,0.1,0.5,1\}$ and terminal gaps $10^{-2},\dots,10^{-6}$; the
conjecture predicts that the midpoint mass falls to zero as the gap shrinks for every $\zeta>0$ and
stays at the ODE value for $\zeta=0$.

\begin{thebibliography}{9}
\bibitem{jain2025attraction} Jain et al.\ (2025), cited in the paper as \texttt{jain2025attraction}.
\end{thebibliography}

\end{document}
